% ============================================================
%  第八章  不完美实验：效应与反事实的定界  （8.1—8.3）
% ============================================================
\chapter{不完美实验：效应与反事实的定界}

\begin{flushright}
\kaishu\small
但愿我发现真理，能像揭穿谬误一样容易。\\[2pt]
——西塞罗（Cicero，公元前 44 年）
\end{flushright}

\section*{序言}
\addcontentsline{toc}{section}{序言}

本章说明，图形模型与反事实模型（\S 3.2 与 \S 7.1）如何结合起来，从不完美实验（imperfect experiments）中获取因果信息——所谓不完美实验，是指偏离随机对照这一理想规程的实验。一种常见的偏离出现在随机临床试验中：受试者并不完全遵从为其分配的处理，从而使因果效应的识别受到损害。当识别的条件不能满足时，我们所能做的最好的事，就是为所关心的量导出界（bounds）——即一个可能取值的范围，它代表我们对数据生成过程的无知，而且不能靠增大样本量来缩小。本章将表明：(i) 这样的界可以用简单的代数方法导出；(ii) 尽管实验并不完美，导出的界仍能就某项政策对整个人群、以及对研究中某个特定个体的影响，给出显著而且有时相当精确的信息；(iii) 先验知识可以被有效地利用，以获得这些影响的贝叶斯估计。

\section{引言}

\subsection{不完美实验与间接实验}

生物学、医学与行为科学中的标准实验研究，总是诉诸随机对照这一手段：受试者被随机分配到各组（或各处理、各项目），不同组参与者之间的均值差，被视为相应项目效力的度量。对这一理想设置的偏离，可能来自未能满足某项实验要求，也可能来自有意放松这些要求。\emph{间接实验}（indirect experiments）是指随机对照或者不可行、或者不可取的研究。在这类实验中，受试者仍被随机分配到各组，但各组成员只是被\emph{鼓励}（而非被强制）参与该组所对应的项目；究竟选择哪个项目，由个体自己决定。

近来，在社会与医学实验中使用严格随机化的做法，受到三方面的质疑。

\begin{enumerate}[label=\arabic*.,leftmargin=2.0em]
\item 完美的控制难以实现，也难以确认。那些假定处理是随机化的研究，可能因未受控制的不完美依从性而受损。例如，对实验药物出现不良反应的受试者，可能决定减少所分配的剂量。又或者，若实验是在检验一种治疗绝症的药物，则怀疑自己被分在对照组的受试者，可能从其他渠道获得该药。这种不完美依从性使实验变成间接实验，并把偏倚引入研究者从数据中得出的结论。除非建立关于依从性的细致模型（Efron and Feldman 1991），这种偏倚无法校正。

\item 不让被分配到某些对照组的受试者获得最佳可得治疗的益处，会带来道德与法律上的后果。例如，在艾滋病研究中，很难为安慰剂方案辩护，因为被分到安慰剂组的病人将被剥夺获得可能救命之治疗的机会（Palca 1989）。

\item 随机化本身的存在，就可能影响参与行为以及受试者的行为（Heckman 1992）。例如，符合条件的候选人一旦发现某所学校刻意对录取标准作随机化，就可能对申请该校心存疑虑。同样，正如 Kramer and Shapiro (1984) 所指出的，药物试验中的受试者参与随机试验的意愿，可能低于参与非实验研究的意愿，即便两种处理同样无害。
\end{enumerate}

总之，研究者开始承认，强制性的随机化可能削弱实验证据的可靠性；而以人为受试者的实验，往往包含——有时也应当包含——自我选择的成分。本章关心的是如何从这样的研究中作出推断：其中受试者对项目有最终选择权，随机化仅限于一个间接的工具（或分配），它仅仅鼓励或劝阻人们参与各个项目。例如，在评估某项培训项目的效力时，可以把资格通知寄给随机抽取的一组学生；或者随机抽取符合条件的候选人，向他们提供参与该项目的奖学金。同样，在药物试验中，受试者可能收到随机选定的关于推荐剂量的建议，但最终的剂量选择仍由受试者根据自身需要决定。

不完美依从性之所以成为问题，是因为仅仅比较处理组与对照组中的比例，可能为“若把该处理统一施加于人群会产生多大效力”给出误导性的估计。例如，如果那些拒绝服药的受试者恰恰是会出现不良反应的人，那么实验可能得出该药比实际更有效的结论。在第三章（\S 3.5，图 3.7(b)）中我们已经说明，此类研究中的处理效力实际上不可识别。也就是说，在没有附加建模假设的情况下，处理效力无法从数据中无偏地估计出来，即便实验中的受试者人数趋于无穷，即便每个受试者的行动与反应都有记录。

本章试图回答的问题是：间接随机化能否提供信息，使我们得以近似评估一个项目的内在价值——例如，当该项目被推广并统一强制施行于人群时所测得的价值。下文的分析将表明，给定一组极少的假设，这样的推断确实可能，只不过其形式是项目或处理之因果效应的界，而非精确的点估计。分析者可以用这些界来保证：某项目的因果效应必定高于某个可测的量，且低于另一个可测的量。

我们的关键假设是：对任何给定的人而言，起鼓励作用的工具会影响此人选择的处理，但不会影响此人对接受到的所选处理会作何反应（参见 \S 7.4.5 中工具变量的定义）。第二个假设在实验研究中总是被采用：受试者对处理的反应彼此独立。除这两个假设之外，我们的模型对“处理反应的倾向”与“处理之间的选择”如何互动，不作任何限制。

\subsection{不依从与意向性治疗}

在处理不完美依从性问题的一种流行而折衷的做法中，研究者进行“意向性治疗”（intent to treat）分析：不管受试者实际是否接受了处理，直接比较对照组与处理组。\footnote{这种做法目前被某些 FDA 机构用于批准新药。}这种分析的结果，度量的是处理\emph{分配}对疾病的作用有多好，而不是我们所期望的“处理\emph{本身}对疾病的作用有多好”。只要实验条件完全模拟该处理最终使用时的条件，基于意向性治疗分析的估计就是有效的。特别是，实验应当模拟受试者接受每种处理的动机。在现场动机比实验动机更强的情况下——通常正是药物获得政府机构批准时的情形——处理效力可能与分配效力相差很大。例如，设想这样一项研究：(a) 该药对人群中很大一部分人有不良副作用；(b) 只有这部分人当中退出处理“臂”（子人群）的成员才康复。意向性治疗分析会把这些康复病例归功于该药，因为它们属于意向性治疗臂；而实际上，这些病例是靠避免该处理才康复的。

处理该问题的另一种做法，是使用基于工具变量公式的校正因子（Angrist et al. 1996）：按此，意向性治疗度量应除以遵从所分配处理的受试者比例。Angrist 等 (1996) 表明，在某些条件下，校正后的公式对“有反应者”（responsive）子人群是有效的——所谓有反应者，是指若被给予不同分配便会改变处理状态的受试者。遗憾的是，这个子人群不可识别；更严重的是，它不能作为涉及整个人群之政策的依据，因为它依赖于工具：在研究中属于有反应者的人，到了现场未必仍是有反应者——现场获得处理的动机与研究所用的动机不同。因此，我们把分析集中在处理中稳定的那个方面，即对依从行为的变化保持不变的方面。

\section{用工具变量界定因果效应}

\subsection{问题的表述：约束优化}

与间接实验相关的基本实验设置如图 8.1 所示，它与图 3.7(b) 和图 5.9 同构。为使讨论集中，我们考虑一项有部分依从性的典型临床试验；不过一般而言，本分析适用于任何这样的研究：其中一个工具变量（定义 7.4.1）鼓励受试者选择某个项目而非另一个。

我们假定 $Z$、$X$、$Y$ 是观测到的二值变量，其中 $Z$ 表示（随机化的）处理分配，$X$ 是实际接受的处理，$Y$ 是观测到的响应。$U$ 项代表影响受试者对处理作何反应的一切因素，无论观测到与否；因此，从 $U$ 到 $Y$ 画一条箭头。从 $U$ 到 $X$ 的箭头表示 $U$ 中的因素也可能影响受试者对处理 $X$ 的选择；这种依赖可能代表介于分配（$Z$）与实际处理（$X$）之间的一个复杂决策过程。

为便于记号，我们让 $z$、$x$、$y$ 分别表示变量 $Z$、$X$、$Y$ 所取的取值，其含义如下：
\begin{flushleft}
\hspace{1.5em}$z\in\{z_0,z_1\}$，$z_1$ 断言处理已被分配（$z_0$ 为其否定）；\\
\hspace{1.5em}$x\in\{x_0,x_1\}$，$x_1$ 断言处理已被施用（$x_0$ 为其否定）；\\
\hspace{1.5em}$y\in\{y_0,y_1\}$，$y_1$ 断言出现正的观测响应（$y_0$ 为其否定）。
\end{flushleft}

$U$ 的定义域仍未具体指定，一般而言，它可能是若干随机变量（离散与连续）之空间的组合。

该图形模型反映两个假设。

\begin{enumerate}[label=\arabic*.,leftmargin=2.0em]
\item 所分配的处理 $Z$ 不直接影响 $Y$，而只通过实际处理 $X$ 影响 $Y$。在实践中，$Z$ 对 $Y$ 可能有的任何直接效应，都可以通过使用安慰剂加以调整。
\item 变量 $Z$ 与 $U$ 边缘独立；这一点由 $Z$ 的随机化保证，随机化排除了 $Z$ 与 $U$ 的共同原因。
\end{enumerate}

\begin{figure}[htbp]
\centering
\begin{tikzpicture}[causal,scale=1.15,every node/.style={transform shape}]
  \node[vtx] (Z) at (0,2.3) {$Z$};
  \node[lat] (U) at (2.4,2.3) {$U$};
  \node[vtx] (X) at (0.7,1.15) {$X$};
  \node[vtx] (Y) at (1.45,0.0) {$Y$};
  \draw[ed] (Z)--(X);
  \draw[edd] (U)--(X);
  \draw[ed] (X)--(Y);
  \draw[edd] (U) to[bend right=22] (Y);
  \node[left=1.2mm,font=\footnotesize,align=right] at (Z.west) {处理\\分配};
  \node[right=1.2mm,font=\footnotesize,align=left] at (U.east) {潜\\因素};
  \node[left=1.2mm,font=\footnotesize,align=right] at (X.west) {接受的\\处理};
  \node[left=1.2mm,font=\footnotesize,align=right] at (Y.west) {观测\\响应};
\end{tikzpicture}
\caption{部分依从的随机临床试验中因果依赖的图形表示。$Z$ 充当工具变量。虚线圆表示未观测变量。}
\label{fig:8.1}
\end{figure}

这些假设对联合分布施加了分解
\begin{equation}
P(y,x,z,u)=P(y\mid x,u)\,P(x\mid z,u)\,P(z)\,P(u), \tag{8.1}
\end{equation}
当然，这个分布无法直接观测，因为 $U$ 未被观测。然而，边缘分布 $P(y,x,z)$、特别是其中的条件分布
\begin{equation}
P(y,x\mid z)=\sum_{u}P(y\mid x,u)\,P(x\mid z,u)\,P(u),\qquad z\in\{z_0,z_1\} \tag{8.2}
\end{equation}
是可观测的，\footnote{在实践中，当然只能观测到 $P(y,x\mid z)$ 的一个有限样本。但我们的任务是识别而非估计，所以我们作大样本假定，把 $P(y,x\mid z)$ 视为已知。}而挑战在于：从这些分布评估处理所导致的 $Y$ 的平均变化。

处理效应由分布 $P(y\mid\doxp(x))$ 支配；利用 \textup{(3.10)} 的截断分解公式，它由下式给出
\begin{equation}
P(y\mid\doxp(x))=\sum_{u}P(y\mid x,u)\,P(u); \tag{8.3}
\end{equation}
这里，因子 $P(y\mid x,u)$ 与 $P(u)$ 就是 \textup{(8.2)} 中的那些因子。因此，如果我们关心的是处理所导致的 $Y$ 的平均变化，就应当计算\emph{平均因果效应}
\[
\ACE(X\to Y)=P(y_1\mid\doxp(x_1))-P(y_1\mid\doxp(x_0))
\]
（Holland 1988），它由下式给出
\begin{equation}
\ACE(X\to Y)=\sum_{u}\bigl[P(y_1\mid x_1,u)-P(y_1\mid x_0,u)\bigr]P(u). \tag{8.4}
\end{equation}

于是，我们的任务就是：给定 \textup{(8.2)} 所表达的观测概率 $P(y,x\mid z_0)$ 与 $P(y,x\mid z_1)$，估计或界定 \textup{(8.4)} 中的表达式。这一任务相当于一个约束优化问题：在 \textup{(8.2)} 的等式约束下，求 \textup{(8.4)} 的最大值与最小值，其中最大化遍及一切满足这些约束的函数
\[
P(u),\ P(y_1\mid x_0,u),\ P(y_1\mid x_1,u),\ P(x_1\mid z_0,u),\ P(x_1\mid z_1,u).
\]

\subsection{规范划分：有限响应变量的演化}

\S 8.2.1 所述的定界问题，可以用数学优化的常规技巧求解。然而，所涉函数的连续性质——以及 $U$ 的未指定定义域——使这种表述不便计算。我们可以换一种做法，利用这样一个观察：$U$ 总可以被一个有限状态变量替代，使得所得模型在关于 $Z$、$X$、$Y$ 的一切观测与操作上都等价（Pearl 1994a）。

考虑因果模型中连接两个二值变量 $Y$ 与 $X$ 的结构方程：
\[
y=f(x,u).
\]
对任意给定的 $u$，$X$ 与 $Y$ 之间的关系必是四个函数之一：
\begin{equation}
f_0:\ y=0,\qquad f_1:\ y=x,\qquad f_2:\ y=\neg x,\qquad f_3:\ y=1. \tag{8.5}
\end{equation}

当 $u$ 沿其定义域变化时，无论这种变化多么复杂，它对模型所能产生的唯一影响，就是在上述四个函数之间切换 $X$ 与 $Y$ 的关系。这就把 $U$ 的定义域划分为四个等价类，如图 8.2 所示——每个类包含对应于同一函数的那些点 $u$。于是我们可以用四状态变量 $R(u)$ 替代 $U$，使每个状态代表四个函数之一。概率 $P(u)$ 会自动转化为概率函数 $P(r)$，$r\in\{0,1,2,3\}$，其值等于分配给与 $r$ 相应之等价类的总权重。这种状态极少的变量 $R$，Balke and Pearl (1994a,b) 称之为“响应”变量，Heckerman and Shachter (1995) 称之为“映射”变量，但“规范划分”（canonical partition）这个名称更具描述性。\footnote{在实验框架中，这种划分可追溯到 Greenland and Robins (1986)，并被 Frangakis and Rubin (2002) 称为“主分层”（Principal Stratification）。在这一框架中（见 \S 7.4.4），$u$ 代表一个实验单元，$R(u)$ 对应于单元 $u$ 对处理 $x$ 的潜在响应。每个单元（例如某个受试者）都具有某种内在的、看似“宿命论式”的响应函数——这一假定曾遭到一些反对（Dawid 2000），理由是对个体处理反应起支配作用的许多因素本质上是不可观测的。$R(u)$ 的等价类表述缓和了这些反对意见（Pearl 2000）：它表明，只要承认这些潜变量的存在（通过方程 $y=f(x,u)$），$R(u)$ 就会从任何复杂的随机潜变量系统中自然地、数学地演化出来。至于那些对后一步也提出量子力学式反对的人（例如 Salmon 1998），不妨把函数关系 $y=f(x,u)$ 视为一个抽象的数学构造，它表示满足 \textup{(8.1)} 与 \textup{(8.2)} 约束的条件概率 $P(y\mid x,u)$ 之集合的极端点（顶点）。}

\begin{figure}[htbp]
\centering
\begin{tikzpicture}[causal,scale=1.0,every node/.style={transform shape}]
  \draw[thick] (0,0) ellipse (3.0 and 1.95);
  \draw[thick] (-0.9,-1.9) .. controls (-1.25,0) .. (-0.9,1.9);
  \draw[thick] (-2.98,0.06) .. controls (-2.1,0.30) and (-1.4,0.06) .. (-0.9,0.06)
               .. controls (-0.2,-0.30) .. (2.98,-0.06);
  \node[align=center,font=\footnotesize] at (-2.0,0.92) {$r_y=1$\\[1pt] $y=x$};
  \node[align=center,font=\footnotesize] at (1.0,0.88) {$r_y=0$\\[1pt] $y=0$};
  \node[align=center,font=\footnotesize] at (-2.0,-0.95) {$r_y=2$\\[1pt] $y=\neg x$};
  \node[align=center,font=\footnotesize] at (1.0,-0.95) {$r_y=3$\\[1pt] $y=1$};
  \draw[ed] (5.0,0.6)--(3.15,0.35);
  \node[font=\footnotesize,anchor=west,align=left] at (5.05,0.65) {$U$ 的定义域};
\end{tikzpicture}
\caption{$U$ 的规范划分：对任意给定函数 $y=f(x,u)$，$U$ 被分为四个等价类，每一类诱导一个从 $X$ 到 $Y$ 的不同函数映射。}
\label{fig:8.2}
\end{figure}

因为 $Z$、$X$、$Y$ 都是二值变量，$U$ 的状态空间划分为 16 个等价类：每个类规定两个函数映射，一个从 $Z$ 到 $X$，另一个从 $X$ 到 $Y$。为描述这些等价类，方便的做法是把每个类视为两个四值变量 $R_x$ 与 $R_y$ 之联合空间中的一个点。变量 $R_x$ 通过如下映射决定受试者的依从行为：
\begin{equation}
x=f_X(z,r_x)=
\begin{cases}
x_0, & r_x=0;\\
x_0, & r_x=1\ \text{且}\ z=z_0,\\
x_1, & r_x=1\ \text{且}\ z=z_1,\\
x_1, & r_x=2\ \text{且}\ z=z_0,\\
x_0, & r_x=2\ \text{且}\ z=z_1,\\
x_1, & r_x=3.
\end{cases} \tag{8.6}
\end{equation}
Imbens and Rubin (1997) 把依从行为为 $r_x=0,1,2,3$ 的受试者分别称为“永不接受者”（never-taker）、“依从者”（complier）、“违抗者”（defier）与“总是接受者”（always-taker）。类似地，变量 $R_y$ 通过如下映射决定受试者的响应行为：
\begin{equation}
y=f_Y(x,r_y)=
\begin{cases}
y_0, & r_y=0;\\
y_0, & r_y=1\ \text{且}\ x=x_0,\\
y_1, & r_y=1\ \text{且}\ x=x_1,\\
y_1, & r_y=2\ \text{且}\ x=x_0,\\
y_0, & r_y=2\ \text{且}\ x=x_1,\\
y_1, & r_y=3.
\end{cases} \tag{8.7}
\end{equation}
Heckerman and Shachter (1995) 把响应行为 $r_y=0,1,2,3$ 分别称为“永不康复”（never-recover）、“受益”（helped）、“受损”（hurt）与“总是康复”（always-recover）。

变量 $R_y$ 的各状态与 \S 7.1（定义 7.1.4）所定义的反事实变量 $Y_{x_1}$ 与 $Y_{x_0}$ 之间的对应关系如下：
\[
Y_{x_1}=
\begin{cases}
y_1, & r_y=1\ \text{或}\ r_y=3,\\
y_0, & \text{其他};
\end{cases}
\qquad
Y_{x_0}=
\begin{cases}
y_1, & r_y=2\ \text{或}\ r_y=3,\\
y_0, & \text{其他}.
\end{cases}
\]

一般而言，响应行为与依从行为未必独立，因此图 8.3 中有双箭头 $R_x\longleftrightarrow R_y$（虚线）。$R_x\times R_y$ 上的联合分布需要 15 个独立参数，而这些参数足以刻画图 8.3 的模型，$P(y,x,z,r_x,r_y)=P(y\mid x,r_y)\,P(x\mid r_x,z)\,P(z)\,P(r_x,r_y)$，因为 $Y$ 与 $X$ 在图中与其父节点处于固定的函数关系。处理的因果效应现在可以直接由 \textup{(8.7)} 得到，即
\begin{equation}
P(y_1\mid\doxp(x_1))=P(r_y=1)+P(r_y=3), \tag{8.8}
\end{equation}
\begin{equation}
P(y_1\mid\doxp(x_0))=P(r_y=2)+P(r_y=3), \tag{8.9}
\end{equation}
以及
\begin{equation}
\ACE(X\to Y)=P(r_y=1)-P(r_y=2). \tag{8.10}
\end{equation}

\begin{figure}[htbp]
\centering
\begin{tikzpicture}[causal,scale=1.15,every node/.style={transform shape}]
  % state boxes
  \node[draw,rounded corners=2pt,minimum width=1.7cm,minimum height=0.8cm,inner sep=1pt] (Rz) at (0.5,3.5) {};
  \foreach \s/\xx in {0/0.05,1/0.95}{\node[font=\scriptsize] at (\xx,3.66) {\s}; \fill (\xx,3.44) circle (1.6pt);}
  \node[font=\small,above=1pt] at (Rz.north) {$R_z$};
  \node[draw,rounded corners=2pt,minimum width=3.0cm,minimum height=0.8cm,inner sep=1pt] (Rx) at (2.1,2.45) {};
  \foreach \s/\xx in {0/1.05,1/1.75,2/2.45,3/3.15}{\node[font=\scriptsize] at (\xx,2.61) {\s}; \fill (\xx,2.39) circle (1.6pt);}
  \node[font=\small,above=1pt] at (Rx.north) {$R_x$};
  \node[draw,rounded corners=2pt,minimum width=3.0cm,minimum height=0.8cm,inner sep=1pt] (Ry) at (3.1,1.35) {};
  \foreach \s/\xx in {0/2.05,1/2.75,2/3.45,3/4.15}{\node[font=\scriptsize] at (\xx,1.51) {\s}; \fill (\xx,1.29) circle (1.6pt);}
  \node[font=\small,below=1pt] at (Ry.south) {$R_y$};
  % observed nodes
  \node[vtx] (Z) at (0.0,2.45) {$Z$};
  \node[vtx] (X) at (1.15,1.15) {$X$};
  \node[vtx] (Y) at (2.05,0.0) {$Y$};
  \draw[ed] (Rz)--(Z);
  \draw[ed] (Z)--(X);
  \draw[ed] (Rx)--(X);
  \draw[ed] (X)--(Y);
  \draw[ed] (Ry)--(Y);
  \draw[bid,dashed] (Rx.east) .. controls (5.1,2.75) and (5.3,1.75) .. (Ry.east);
\end{tikzpicture}
\caption{与图 8.1 等价的结构，但使用有限状态的响应变量 $R_z$、$R_x$ 与 $R_y$。}
\label{fig:8.3}
\end{figure}

\subsection{线性规划表述}

通过阐明 $P(y,x\mid z)$ 的参数与 $P(r_x,r_y)$ 的参数之间的关系，我们得到在给定 $P(y,x\mid z)$ 下最小化或最大化 $\ACE(X\to Y)$ 所需的一组线性约束。

观测变量上的条件分布 $P(y,x\mid z)$ 完全由八个参数确定，记作：
\begin{align*}
p_{00.0}&=P(y_0,x_0\mid z_0), & p_{00.1}&=P(y_0,x_0\mid z_1),\\
p_{01.0}&=P(y_0,x_1\mid z_0), & p_{01.1}&=P(y_0,x_1\mid z_1),\\
p_{10.0}&=P(y_1,x_0\mid z_0), & p_{10.1}&=P(y_1,x_0\mid z_1),\\
p_{11.0}&=P(y_1,x_1\mid z_0), & p_{11.1}&=P(y_1,x_1\mid z_1).
\end{align*}
概率约束
\[
\sum_{n=00}^{11}p_{n.0}=1
\qquad\text{与}\qquad
\sum_{n=00}^{11}p_{n.1}=1
\]
进一步意味着 $\vec{p}=(p_{00.0},\dots,p_{11.1})$ 可以由 6 维空间中的一个点确定。这个空间记作 $\mathcal{P}$。

联合概率 $P(r_x,r_y)$ 有 16 个参数：
\begin{equation}
q_{jk}\triangleq P(r_x=j,\ r_y=k), \tag{8.11}
\end{equation}
其中 $j,k\in\{0,1,2,3\}$。概率约束
\[
\sum_{j=0}^{3}\sum_{k=0}^{3}q_{jk}=1
\]
意味着 $\vec{q}$ 确定 15 维空间中的一个点。这个空间记作 $\mathcal{Q}$。

式 \textup{(8.10)} 现在可以改写为 $\mathcal{Q}$ 参数的一个线性组合：
\begin{equation}
\ACE(X\to Y)=q_{01}+q_{11}+q_{21}+q_{31}-q_{02}-q_{12}-q_{22}-q_{32}. \tag{8.12}
\end{equation}

应用 \textup{(8.6)} 与 \textup{(8.7)}，我们可以写出从 $\mathcal{Q}$ 中的一点到 $\mathcal{P}$ 中的一点的线性变换：
\begin{align*}
p_{00.0}&=q_{00}+q_{01}+q_{10}+q_{11}, & p_{00.1}&=q_{00}+q_{01}+q_{20}+q_{21},\\
p_{01.0}&=q_{20}+q_{22}+q_{30}+q_{32}, & p_{01.1}&=q_{10}+q_{12}+q_{30}+q_{32},\\
p_{10.0}&=q_{02}+q_{03}+q_{12}+q_{13}, & p_{10.1}&=q_{02}+q_{03}+q_{22}+q_{23},\\
p_{11.0}&=q_{21}+q_{23}+q_{31}+q_{33}, & p_{11.1}&=q_{11}+q_{13}+q_{31}+q_{33},
\end{align*}
它也可以写成矩阵形式 $\vec{p}=\vec{R}\vec{q}$。

给定 $\mathcal{P}$ 空间中的一点 $\vec{p}$，$\ACE(X\to Y)$ 的严格下界可以通过求解如下线性规划问题确定。
\begin{equation}
\begin{aligned}
&\text{最小化}\quad q_{01}+q_{11}+q_{21}+q_{31}-q_{02}-q_{12}-q_{22}-q_{32}\\
&\text{约束于：}\quad \sum_{j=0}^{3}\sum_{k=0}^{3}q_{jk}=1,\qquad \vec{R}\vec{q}=\vec{p},\\
&\hspace{4.2em} q_{jk}\ge 0\quad\text{对}\ j,k\in\{0,1,2,3\}.
\end{aligned} \tag{8.13}
\end{equation}

对于这种规模的问题，有现成程序可以导出该优化问题解的符号表达式（Balke 1995），从而得到处理效应的如下下界：
\begin{equation}
\ACE(X\to Y)\ \ge\ \max\left\{
\begin{array}{l}
p_{11.1}+p_{00.0}-1\\
p_{11.0}+p_{00.1}-1\\
p_{11.0}-p_{11.1}-p_{10.1}-p_{01.0}-p_{10.0}-p_{01.1}-p_{10.1}\\
p_{11.1}-p_{11.0}-p_{10.0}-p_{01.1}-p_{10.1}-p_{01.0}-p_{10.0}\\
-p_{01.1}-p_{10.1}\\
-p_{01.0}-p_{10.0}\\
p_{00.1}-p_{01.1}-p_{10.1}-p_{01.0}-p_{00.0}\\
p_{00.0}-p_{01.0}-p_{10.0}-p_{01.1}-p_{00.1}
\end{array}\right. \tag{8.14a}
\end{equation}
类似地，上界由下式给出
\begin{equation}
\ACE(X\to Y)\ \le\ \min\left\{
\begin{array}{l}
1-p_{01.1}-p_{10.0}\\
1-p_{01.0}-p_{10.1}\\
-p_{01.0}+p_{01.1}+p_{00.1}+p_{11.0}+p_{00.0}\\
-p_{01.1}+p_{11.1}+p_{00.1}+p_{01.0}+p_{00.0}\\
p_{11.1}+p_{00.1}\\
p_{11.0}+p_{00.0}\\
-p_{10.1}+p_{11.1}+p_{00.1}+p_{11.0}+p_{10.0}\\
-p_{10.0}+p_{11.0}+p_{00.0}+p_{11.1}+p_{10.1}
\end{array}\right. \tag{8.14b}
\end{equation}

我们也可以在同一组线性约束下，分别导出 \textup{(8.8)} 与 \textup{(8.9)} 的界，得到：
\begin{equation}
\begin{aligned}
P(y_1\mid\doxp(x_0))&\ \ge\ \max\left\{
\begin{array}{l}
p_{10.0}+p_{11.0}-p_{00.1}-p_{11.1}\\
p_{10.1}\\
p_{10.0}\\
p_{01.0}+p_{10.0}-p_{00.1}-p_{01.1}
\end{array}\right\},\\
P(y_1\mid\doxp(x_0))&\ \le\ \min\left\{
\begin{array}{l}
p_{01.0}+p_{10.0}+p_{10.1}+p_{11.1}\\
1-p_{00.1}\\
1-p_{00.0}\\
p_{10.0}+p_{11.0}+p_{01.1}+p_{10.1}
\end{array}\right\},
\end{aligned} \tag{8.15}
\end{equation}
\begin{equation}
\begin{aligned}
P(y_1\mid\doxp(x_1))&\ \ge\ \max\left\{
\begin{array}{l}
p_{11.0}\\
p_{11.1}\\
-p_{00.0}-p_{01.0}+p_{00.1}+p_{11.1}\\
-p_{01.0}-p_{10.0}+p_{10.1}+p_{11.1}
\end{array}\right\},\\
P(y_1\mid\doxp(x_1))&\ \le\ \min\left\{
\begin{array}{l}
1-p_{01.0}\\
1-p_{01.1}\\
p_{00.0}+p_{11.0}+p_{10.1}+p_{11.1}\\
p_{10.0}+p_{11.0}+p_{00.1}+p_{11.1}
\end{array}\right\}.
\end{aligned} \tag{8.16}
\end{equation}
这些表达式给出了所寻求之量在无假设条件下最紧的界。\footnote{“假设透明”（assumption-transparent）也许是更好的说法；我们不对决定受试者依从性的因素作任何假设，但我们依赖 (i) 随机分配与 (ii) 无副作用这两个假设，如图中所示（例如图 8.1）。}

\subsection{自然界}

$\ACE(X\to Y)$（式 \textup{(8.4)}）的表达式可以用两个简单公式定界，它们各由 \textup{(8.14a)} 与 \textup{(8.14b)} 中的前两项构成（Robins 1989; Manski 1990; Pearl 1994a）：
\begin{equation}
\begin{aligned}
\ACE(X\to Y)&\ \ge\ P(y_1\mid z_1)-P(y_1\mid z_0)-P(y_1,x_0\mid z_1)-P(y_0,x_1\mid z_0),\\
\ACE(X\to Y)&\ \le\ P(y_1\mid z_1)-P(y_1\mid z_0)+P(y_0,x_0\mid z_1)+P(y_1,x_1\mid z_0).
\end{aligned} \tag{8.17}
\end{equation}

由于形式简单、适用范围广，\textup{(8.17)} 给出的界被称为\emph{自然界}（natural bounds）（Balke and Pearl 1997）。自然界保证：实际处理的因果效应不可能比鼓励的因果效应（$P(y_1\mid z_1)-P(y_1\mid z_0)$）小出两个可测量 $P(y_1,x_0\mid z_1)+P(y_0,x_1\mid z_0)$ 之和以上；它们也保证：处理的因果效应不可能比鼓励的因果效应大出另外两个可测量 $P(y_0,x_0\mid z_1)+P(y_1,x_1\mid z_0)$ 之和以上。自然界的宽度，不出所料，由不依从率给出：$P(x_1\mid z_0)+P(x_0\mid z_1)$。

不过，\textup{(8.14ab)} 中锐界（sharp bounds）的宽度可能窄得多。Balke (1995) 与 Pearl (1995b) 表明，即便在不依从率达 50\% 的条件下，这些界也可能收缩为一个点，从而允许对 $\ACE(X\to Y)$ 作一致估计。只要 (a) 遵从分配 $z_0$ 的受试者比例与遵从 $z_1$ 的相同，且 (b) 在至少一个处理臂 $x$ 中 $Y$ 与 $Z$ 完全相关（见表 8.1，\S 8.5），这种情况就会发生。

\textup{(8.14ab)} 的锐界虽然比 \textup{(8.17)} 的自然界复杂，但一旦掌握了 $P(y,x\mid z)$ 八个格子中的频数数据，就不难计算。还可以证明（Balke 1995）：当我们能够安全地假定没有受试者是“反其道者”（contrarian）——换言之，没有受试者会一贯地选择与所分配处理相反的处理臂——自然界就是最优的。

注意，如果响应 $Y$ 是连续的，就可以把 $y_1$ 与 $y_0$ 分别与二值事件 $Y>t$ 和 $Y\le t$ 相联系，并让 $t$ 在 $Y$ 的范围内连续变化。此时 \textup{(8.15)} 与 \textup{(8.16)} 将给出处理效应 $P(Y\le t\mid\doxp(x))$ 之整个分布的界。

\subsection{处理对受处理者的效应（ETT）}

大量文献假定 $\ACE(X\to Y)$ 是所关心的参数，因为 $\ACE(X\to Y)$ 预测了把处理统一（或随机）施加于人群时的影响。然而，如果政策制定者并不关心引入新的处理政策，而是关心在现行激励机制下维持还是终止一个既有项目，那么所关心的参数就应当度量处理对受处理者的影响，即受处理受试者的平均响应与这些\emph{同一批}受试者在未受处理时的平均响应之差（Heckman 1992）。该参数的恰当公式是
\begin{equation}
\begin{aligned}
\ETT(X\to Y)&=P(Y_{x_1}=y_1\mid x_1)-P(Y_{x_0}=y_1\mid x_1)\\
&=\sum_{u}\bigl[P(y_1\mid x_1,u)-P(y_1\mid x_0,u)\bigr]P(u\mid x_1),
\end{aligned} \tag{8.18}
\end{equation}
它与 \textup{(8.4)} 相似，只是把对 $u$ 的期望换成了给定 $X=x_1$ 的条件期望。

对 $\ETT(X\to Y)$ 的分析表明，在无侵入（no intrusion）的条件下（即 $P(x_1\mid z_0)=0$，如大多数临床试验那样），$\ETT(X\to Y)$ 可以被精确识别（Bloom 1984; Heckman and Robb 1986; Angrist and Imbens 1991）。一般情形下支配 $\ETT(X\to Y)$ 的界可以用类似手段得到（Pearl 1995b），其结果为
\begin{equation}
\begin{aligned}
\ETT(X\to Y)&\ \ge\ \frac{P(y_1\mid z_1)-P(y_1\mid z_0)}{P(x_1)/P(z_1)}-\frac{P(y_0,x_1\mid z_0)}{P(x_1)},\\
\ETT(X\to Y)&\ \le\ \frac{P(y_1\mid z_1)-P(y_1\mid z_0)}{P(x_1)/P(z_1)}+\frac{P(y_1,x_1\mid z_0)}{P(x_1)}.
\end{aligned} \tag{8.19}
\end{equation}

锐界见 Balke (1995, p.~113)。显然，在处理只能由（因分配而）受鼓励者获得的情形下，我们有 $P(x_1\mid z_0)=0$，且
\begin{equation}
\ETT(X\to Y)=\frac{P(y_1\mid z_1)-P(y_1\mid z_0)}{P(x_1\mid z_1)}. \tag{8.20}
\end{equation}
与 $\ACE(X\to Y)$ 不同，$\ETT(X\to Y)$ 不是处理的内在属性，因为它随鼓励工具而变化。因此，它的意义在于那些希望评估既有项目对其当前参与者之效力的研究。

\subsection{例：消胆胺的效应}

为用例子说明 $\ACE(X\to Y)$ 的界如何能提供有关因果效应的有意义信息，我们考虑脂质研究诊所冠心病一级预防试验的数据（Program 1984）。Efron and Feldman (1991) 分析了其中一部分（涵盖 337 名受试者），这也是本例的焦点。受试者被随机分入两个人数大致相等的治疗组；一组中所有受试者都被处方消胆胺（$z_1$），另一组受试者被处方安慰剂（$z_0$）。在数年的治疗过程中，每名受试者的胆固醇水平被多次测量，这些测量的平均值被用作治疗后胆固醇水平（连续变量 CF）。每名受试者的依从性由追踪其所服用的处方剂量（一个连续量）来确定。

为把 \textup{(8.17)} 的界应用于本研究的数据，先用阈值把连续数据变换为代表处理分配（$Z$）、接受的处理（$X$）与处理响应（$Y$）的二值变量。剂量消耗的阈值选在大致为最小与最大消耗量的中点处；胆固醇水平下降的阈值定为 28 个单位。经过这一“阈值化”程序后，数据样本给出如下八个概率：\footnote{我们作大样本假定，以样本频率代表 $P(y,x\mid z)$。为考虑样本变异性，一切界都应当补充置信区间与显著性水平，如同对照实验的传统分析那样。\S 8.5.1 用吉布斯抽样评估样本变异性。}
\begin{align*}
P(y_0,x_0\mid z_0)&=0.919, & P(y_0,x_0\mid z_1)&=0.315,\\
P(y_0,x_1\mid z_0)&=0.000, & P(y_0,x_1\mid z_1)&=0.139,\\
P(y_1,x_0\mid z_0)&=0.081, & P(y_1,x_0\mid z_1)&=0.073,\\
P(y_1,x_1\mid z_0)&=0.000, & P(y_1,x_1\mid z_1)&=0.473.
\end{align*}
这些数据代表的依从率为
\[
P(x_1\mid z_1)=0.139+0.473=0.61,
\]
均值差（用 $P(z_1)=0.50$）为
\[
P(y_1\mid x_1)-P(y_1\mid x_0)=\frac{0.473}{0.473+0.139}-\frac{0.073+0.081}{0.919+0.315+0.081+0.073}=0.662,
\]
而鼓励效应（意向性治疗）为
\[
P(y_1\mid z_1)-P(y_1\mid z_0)=0.073+0.473-0.081=0.465.
\]
根据 \textup{(8.17)}，$\ACE(X\to Y)$ 可以被界定为
\[
\begin{aligned}
\ACE(X\to Y)&\ \ge\ 0.465-0.073-0.000=0.392,\\
\ACE(X\to Y)&\ \le\ 0.465+0.315+0.000=0.780.
\end{aligned}
\]
这些界的信息量大得惊人：尽管 38.8\% 的受试者偏离了其处理规程，实验者仍可以明确断言，当统一施加于人群时，该处理保证至少使胆固醇水平下降 28 点以上的概率提高 39.2\%。

处理“对受处理者”的影响同样发人深省。用式 \textup{(8.20)}，$\ETT(X\to Y)$ 可以被精确计算（因为 $P(x_1\mid z_0)=0$）：
\[
\ETT(X\to Y)=\frac{0.465}{0.610}=0.762.
\]
换言之，那些留在项目中的受试者，其境况比未受处理时好得多：在这些人当中，至少 76.2\% 的人可以归功于该处理而使胆固醇水平下降至少 28 个单位。

\section{反事实与法律责任}

在一些法律案件中，评估反事实概率可能颇具启发性——原告声称被告的行为应对原告的不幸负责。若对反事实处理不当，很容易作出不当裁决（Robins and Greenland 1989）。考虑下面这个假想的、虚构的案例研究，它由 Balke and Pearl (1994a) 专门构造，以突出因果效应与因果归因之间的差异。

PeptAid（一种抗酸药）的营销商向加州 Stress 城 10\% 的家庭随机邮寄了产品样品。在一项后续研究中，研究者确定每个人是否收到了 PeptAid 样品、是否服用了 PeptAid、以及在此后一个月内是否发生消化性溃疡。

这一情景的因果结构与图 8.1 给出的部分依从模型完全相同，其中 $z_1$ 断言从营销商处收到了 PeptAid，$x_1$ 断言服用了 PeptAid，$y_1$ 断言发生了消化性溃疡。数据显示如下分布：
\begin{align*}
P(y_0,x_0\mid z_0)&=0.32, & P(y_0,x_0\mid z_1)&=0.02,\\
P(y_0,x_1\mid z_0)&=0.32, & P(y_0,x_1\mid z_1)&=0.17,\\
P(y_1,x_0\mid z_0)&=0.04, & P(y_1,x_0\mid z_1)&=0.67,\\
P(y_1,x_1\mid z_0)&=0.32, & P(y_1,x_1\mid z_1)&=0.14.
\end{align*}
这些数据表明，服用 PeptAid 者与发生消化性溃疡者之间存在高度相关：
\[
P(y_1\mid x_1)=0.50,\qquad P(y_1\mid x_0)=0.26.
\]
此外，意向性治疗分析表明，收到 PeptAid 样品的人发生消化性溃疡的机会高出 45\%：
\[
P(y_1\mid z_1)=0.81,\qquad P(y_1\mid z_0)=0.36.
\]

原告（Smith 先生）得知这项研究后，同时对营销公司与 PeptAid 生产商提起诉讼。原告律师针对生产商辩称，服用 PeptAid 触发了其当事人的溃疡及由此产生的医疗费用。同样，原告律师针对营销商辩称，若营销商未分发产品样品，其当事人本不会发生溃疡。

代表 PeptAid 制造商与营销商双方的辩护律师反驳说，PeptAid 服用与溃疡之间的高度相关应归因于一个共同因素，即溃疡前的不适。有胃肠不适的人，既更可能使用 PeptAid，也更可能发生胃溃疡。为支持其当事人的主张，辩护律师提交了对数据的专家分析，表明平均而言，服用 PeptAid 实际上使个体发生溃疡的机会至少降低 15\%。

确实，应用 \textup{(8.14a,b)} 可得到 PeptAid 服用对消化性溃疡之平均因果效应的如下界：
\[
-0.23\ \le\ \ACE(X\to Y)\ \le\ -0.15;
\]
这证明 PeptAid 对整个人群是有益的。

不过，原告律师强调了整个人群的平均因果效应与一个子人群之平均因果效应之间的区别，该子人群由那些像他当事人一样收到 PeptAid 样品、服用之、并随后发生溃疡的个体组成。对人群数据的分析表明，若未分发 PeptAid，Smith 先生发生溃疡的机会至多为 7\%——无论存在何种混杂因子（如溃疡前疼痛）。同样，若 Smith 先生未服用 PeptAid，他发生溃疡的机会至多为 7\%。

不利于营销商的有害统计量，来自对如下反事实概率之界的评估：在原告确实收到了 PeptAid 样品、服用了 PeptAid、并发生了消化性溃疡的条件下，若他未曾收到 PeptAid 样品，他本会发生消化性溃疡的概率。该概率可以用参数 $q_{13}$、$q_{31}$ 与 $q_{33}$ 表示为
\[
P(Y_{z_0}=y_1\mid z_1,x_1,y_1)=\frac{P(r_z=1)(q_{13}+q_{31}+q_{33})}{P(y_1,x_1,z_1)},
\]
因为只有组合 $\{r_x=1,r_y=3\}$、$\{r_x=3,r_y=1\}$ 与 $\{r_x=3,r_y=3\}$ 满足联合事件 $\{X=x_1,\ Y=y_1,\ Y_{z_0}=y_1\}$（见 \textup{(8.6)}、\textup{(8.7)} 与 \textup{(8.11)}）。因此，
\[
P(Y_{z_0}=y_1\mid z_1,x_1,y_1)=\frac{q_{13}+q_{31}+q_{33}}{P(y_1,x_1\mid z_1)}.
\]
该表达式对 $q$ 参数是线性的，可以用线性规划定界，得到
\begin{equation*}
\begin{aligned}
P(Y_{z_0}=y_1\mid z_1,x_1,y_1)&\ \ge\ \frac{1}{p_{11.1}}\max\left\{
\begin{array}{l}
0\\
p_{11.1}-p_{00.0}\\
p_{11.0}-p_{00.1}-p_{10.1}\\
p_{10.0}-p_{01.1}-p_{10.1}
\end{array}\right\},\\
P(Y_{z_0}=y_1\mid z_1,x_1,y_1)&\ \le\ \frac{1}{p_{11.1}}\min\left\{
\begin{array}{l}
p_{11.1}\\
p_{10.0}+p_{11.0}\\
1-p_{00.0}-p_{10.1}
\end{array}\right\}.
\end{aligned}
\end{equation*}
类似地，不利于 PeptAid 生产商的有害证据，来自评估反事实概率
\[
P(Y_{x_0}=y_1\mid y_1,x_1,z_1)=\frac{q_{13}+q_{33}}{p_{11.1}}
\]
的界。若对分子作最小化与最大化（在 \textup{(8.13)} 的约束下），我们得到
\begin{equation*}
\begin{aligned}
P(Y_{x_0}=y_1\mid y_1,x_1,z_1)&\ \ge\ \frac{1}{p_{11.1}}\max\left\{
\begin{array}{l}
0\\
p_{11.1}-p_{00.0}-p_{11.0}\\
p_{10.0}-p_{01.1}-p_{10.1}
\end{array}\right\},\\
P(Y_{x_0}=y_1\mid y_1,x_1,z_1)&\ \le\ \frac{1}{p_{11.1}}\min\left\{
\begin{array}{l}
p_{11.1}\\
p_{10.0}+p_{11.0}\\
1-p_{00.0}-p_{10.1}
\end{array}\right\}.
\end{aligned}
\end{equation*}
把观测分布 $P(y,x\mid z)$ 代入这些公式，得到如下界：
\[
0.93\ \le\ P(Y_{z_0}=y_0\mid z_1,x_1,y_1)\ \le\ 1.00,
\]
\[
0.93\ \le\ P(Y_{x_0}=y_0\mid z_1,x_1,y_1)\ \le\ 1.00.
\]
因此，原告这一类人中至少 93\% 若未被鼓励服用 PeptAid（$z_0$）、或者同样地若未服用 PeptAid（$x_0$），本不会发生溃疡。这为原告的主张提供了非常有力的支持：他确实受到了营销商与生产商之行为及产品的伤害。

第九章将继续分析具体事件中的因果归因，并将建立这样的条件：在这些条件下，正确归因的概率可以从实验数据与非实验数据中识别出来。
